\documentclass[11pt]{article}
\usepackage[a4paper,margin=26mm]{geometry}
\usepackage{amsmath,amssymb}
\usepackage[hidelinks]{hyperref}
\setlength{\emergencystretch}{3em}
\title{Conjecture 00000001491: the printed map collapses onto a line}
\author{lizaixi01, with OpenAI Codex assistance}
\date{23 October 2026}
\begin{document}
\maketitle
\section{The explicit formula matters}
Both source languages explicitly define
\[
 H_{a,b}(x,y)=(1-ax^2+y,\ b).
\]
\textbf{The second coordinate here is the constant $b$, not $bx$.}
The usual H\'enon map in H\'enon's original abstract~\cite{Henon}
has second coordinate $bx$. Our result concerns exactly the formula
printed in this problem; it does not assert a dimension theorem for
that different, usual H\'enon map.

For the printed map, every image lies on the horizontal line
$L_b=\{(x,y):y=b\}$. Consequently every positive-time iterate,
every actual forward $\omega$-limit set, and the closure of the union
of all actual point $\omega$-limit sets lie on $L_b$. Their Hausdorff
dimensions are at most one. Any compact fully invariant set
$H_{a,b}(A)=A$ also has dimension at most one. These statements are
proved for all real $a,b$, then instantiated at the exact printed
parameters $a=7/5=1.4$, $b=3/10=0.3$.

Thus the required lower bound $1.26<d_H$ cannot hold for the
printed map's actual limiting sets or a fully invariant compact
attractor. This necessary lower-bound clause already refutes the
conjunction. We do not calculate the usual H\'enon attractor's
dimension or assert independent results about box dimension or
the Kaplan--Yorke formula. No classification of all alternative
Milnor-style attractor definitions is claimed.

\section{The actual Euclidean plane and horizontal line}
Lean uses the actual Euclidean space
$E=\operatorname{EuclideanSpace}(\mathbb R,\operatorname{Fin}2)$,
with its Euclidean metric, and the actual continuous polynomial map
given above. The line $L_b\subset E$ is closed because the second
coordinate projection is continuous. For every $v\in E$,
\[
 H_{a,b}(v)\in L_b,\qquad
 H_{a,b}^{\,n}(v)\in L_b\quad(n\ge1).
\]
The second statement follows from the actual iterate successor
identity, not an assumed finite-state or line-restricted model.

The genuine line parametrization
\[
 \iota_b:\mathbb R\longrightarrow E,\qquad t\longmapsto(t,b)
\]
is an actual isometry onto $L_b$. Lean proves the distance identity
using the Euclidean squared-coordinate formula, proves
$L_b=\iota_b(\mathbb R)$, and applies isometric invariance of
Hausdorff dimension and $\dim_H\mathbb R=1$. Therefore
\[
 \dim_H L_b=1,\qquad A\subseteq L_b\Rightarrow\dim_H A\le1.
\]
No max-norm substitution or unproved metric conversion is used.

\newpage
\section{Actual forward limit sets}
For any actual initial set $S\subset E$, define
\[
 \Omega_{a,b}(S)=\operatorname{omegaLimit}
 \bigl(\operatorname{atTop},\ (n,v)\mapsto H_{a,b}^{\,n}(v),\ S\bigr).
\]
This is Mathlib's actual filter-based $\omega$-limit set: the
intersection, over $U\in\operatorname{atTop}$, of the closures of
$\{H_{a,b}^{\,n}(v):n\in U,v\in S\}$.
Take the actual filter member $U=\{n:n\ge1\}$. All these images
lie in the closed line $L_b$, so their closure does too. Hence
\[
 \Omega_{a,b}(S)\subseteq L_b,\qquad
 \dim_H\Omega_{a,b}(S)\le1\quad\text{for every }S\subset E.
\]
In particular this applies to every single orbit $S=\{v\}$.
The separately defined actual global limiting object
\[
 K_{a,b}=\overline{\bigcup_{v\in E}\Omega_{a,b}(\{v\})}
\]
also lies in $L_b$ by closedness, and has Hausdorff dimension at
most one. No arbitrary set is simply named an attractor in place
of these actual dynamical objects.

\section{Full invariance and the real dimension bound}
For every actual set $A\subset E$ with $H_{a,b}(A)=A$, each
$v\in A$ is an actual image $H_{a,b}(w)$ for some $w\in A$.
Thus $A\subseteq L_b$. In particular there is no compact fully
invariant set with dimension greater than $63/50=1.26$.
This uses full invariance; merely forward invariance
$H_{a,b}(A)\subseteq A$ would not suffice, since a thick trapping
region may contain points off the line. We do not use that weaker
condition as a substitute for an attractor.

Mathlib's \texttt{dimH} takes values in the extended nonnegative
reals. Lean first proves the actual bound $\dim_H A\le1$ for every
line subset, and separately proves $\dim_H A<\infty$.
The real-value bridge then gives
$\operatorname{toReal}(\dim_H A)\le1$. The final theorem
\begin{center}\small\texttt{source\_parameters\_real\_dimensions}\end{center}
explicitly gives this real bound for every actual $\Omega_{7/5,3/10}(S)$ and for
$K_{7/5,3/10}$, and denies existence of a compact fully invariant
$A$ with real dimension above $63/50$. There is no use of an
infinite-dimension conversion default.

\section{Verification}
Lean 4.33.0 and Mathlib commit
\begin{center}\small\texttt{db584cd6d46c92f209a44c0f1c829460d327499d}\end{center}
are pinned. The package includes a new root build, transitive axiom
audits of all 21 theorems and full official
\texttt{leanchecker --fresh Main} replay. No \texttt{sorry},
\texttt{native\_decide} or additional axiom is used. OpenAI Codex
assisted; a separate supervising Codex agent reviewed source,
report and PDF. This is AI review, not human certification.
Organizer acceptance is pending.
\begin{thebibliography}{1}
\bibitem{Henon} M.~H\'enon, \emph{A two-dimensional mapping with a
strange attractor}, Commun. Math. Phys. \textbf{50} (1976), 69--77.
\url{https://doi.org/10.1007/BF01608556}.
\end{thebibliography}
\end{document}
